
\SGARefTarget{sga7.sga7i.eix.n11-8}{11.8}. As in \SGARefLink{sga7.sga7i.eix.n8-3}{8.3}, regard the separable closure $\overline K$ of $K$ as the
inductive limit of its finite subextensions $K'$. For every such $K'$, let
$A[K']$ denote the N\'{e}ron model of

\begin{equation*}
A_{K'}=A_K\otimes_K K'.
\end{equation*}
It gives a finite \'{e}tale group scheme

\begin{equation}\tag{\SGARefTarget{sga7.sga7i.eix.n11-8-1}{11.8.1}}
\Phi_o[K']=A[K']_o/A[K']_o^o
\end{equation}
of connected components, defined over the residue field of the normalization
$S[K']$ of $S$ in $K'$. For simplicity, suppose that the residue field of
$S$ is separably closed, so that the same is true of the residue fields of
the $S[K']$. The groups (\SGARefLink{sga7.sga7i.eix.n11-8-1}{11.8.1}) may then be identified with ordinary finite
groups $\Phi[K']$, and it is immediate that, as $K'$ varies, these form an
inductive system, which we now propose to study. We shall be interested in
particular in the group

\begin{equation}\tag{\SGARefTarget{sga7.sga7i.eix.n11-8-2}{11.8.2}}
\Psi=\varinjlim_{K'}\Phi[K'].
\end{equation}
We know from \SGARefLink{sga7.sga7i.eix.n3-6}{3.6} that, for $K'$ sufficiently large, say
$K'\supset K'_o$, $A[K']$ has semistable reduction. Hence for
$K'\supset K'_o$ the transition morphisms are injective (\SGARefLink{sga7.sga7i.eix.n3-3}{3.3}). Moreover,
with the notation of \SGARefLink{sga7.sga7i.eix.n5-5}{5.5}, retained in the present section, there is for
$K'\supset K'_o$ a canonical isomorphism

\begin{equation}\tag{\SGARefTarget{sga7.sga7i.eix.n11-8-3}{11.8.3}}
\Phi[K'](\ell)\simeq
\ker\bigl(\widetilde u[K']_\ell\otimes D_\ell:
 M_\ell\otimes D_\ell\longrightarrow
 \check M'_\ell\otimes D_\ell\bigr),
\end{equation}
coming from (\SGARefLink{sga7.sga7i.eix.n11-6-2}{11.6.2}) applied to $A[K']$. One also verifies immediately that
the corresponding immersions

\begin{equation*}
\Phi[K'](\ell)\lhook\joinrel\longrightarrow M_\ell\otimes D_\ell
\end{equation*}
are compatible with the transition morphisms as $K'$ varies. Passing to the
limit therefore gives a canonical monomorphism

\begin{equation}\tag{\SGARefTarget{sga7.sga7i.eix.n11-8-4}{11.8.4}}
\Psi(\ell)\mathrel{\mathop{=}^{\mathrm{def}}}
 \varinjlim_{K'}\Phi[K'](\ell)
 \lhook\joinrel\longrightarrow M_\ell\otimes D_\ell
\end{equation}
from the $\ell$-primary component of $\Psi$; summing over all $\ell$ gives

\begin{equation}\tag{\SGARefTarget{sga7.sga7i.eix.n11-8-5}{11.8.5}}
\Psi\mathrel{\mathop{=}^{\mathrm{def}}}
 \varinjlim_{K'}\Phi[K']
 \longrightarrow M\otimes\mathbb Q/\mathbb Z.
\end{equation}

\medskip
\noindent
\underline{Theorem} \SGARefTarget{sga7.sga7i.eix.n11-9}{11.9}. \underline{The canonical homomorphism (\SGARefLink{sga7.sga7i.eix.n11-8-5}{11.8.5})
is an isomorphism}.

We already know that it is a monomorphism. To prove that it is an
isomorphism, we must use the characterization (\SGARefLink{sga7.sga7i.eix.n11-8-3}{11.8.3}) of the image of
$\Phi[K'](\ell)$ and the behaviour of $\widetilde u[K']$ as $K'$ varies.
That behaviour is given by (\SGARefLink{sga7.sga7i.eix.n10-3-4}{10.3.4}).

We may reduce to the case in which $S$ is strictly local. To prove \SGARefLink{sga7.sga7i.eix.n11-9}{11.9} we
are then reduced to the following familiar assertion, applied to $S[K'_o]$:
if $S$ is a strictly local trait with fraction field $K$, then for every
integer $n\geq1$ there is a separable extension $K'$ of $K$ whose degree is
divisible by $n$.
\footnote{The printed source appends ``cf. I'' followed by a blank internal
reference. No target is supplied, so none is invented here.}

\medskip
\noindent
\underline{Corollary} \SGARefTarget{sga7.sga7i.eix.n11-10}{11.10}. \underline{Under the hypotheses ($S$
complete) and with the notation of \SGARefLink{sga7.sga7i.eix.n8-3}{8.3}, the ind-finite group scheme over
$K$ of connected components of the group scheme $A_K^{\flat}$,
namely $A_K^{\flat}/{A_K^{\flat}}^o$, is canonically isomorphic to}
$\underline M_K\otimes\mathbb Q/\mathbb Z$, \underline{where
$\underline M_K$ is the twisted constant group over $K$ defined in
(\SGARefLink{sga7.sga7i.eix.n8-1-8}{8.1.8})}.

This follows immediately from (8.2.4) and \SGARefLink{sga7.sga7i.eix.n11-9}{11.9}, in view of the definition
of $\underline M_K$.

\medskip
\noindent
\underline{Remarks} \SGARefTarget{sga7.sga7i.eix.n11-11}{11.11}. a) If one restricts to the $\ell$-primary
components for $\ell\neq p$, Theorem \SGARefLink{sga7.sga7i.eix.n11-9}{11.9} is immediate in terms of \SGARefLink{sga7.sga7i.eix.n11-2}{11.2} and
does not require the more delicate considerations of \SGARefLink{sga7.sga7i.eix.n11-6}{11.6} (via \SGARefLink{sga7.sga7i.eix.n9-4}{9.4}).
Likewise, if $K$ is assumed to be \underline{of characteristic zero}, \SGARefLink{sga7.sga7i.eix.n11-9}{11.9}
can be proved without using (\SGARefLink{sga7.sga7i.eix.n11-8-3}{11.8.3}): it is immediate that the image of
${}_n\Phi[K'](\ell)$ in

\begin{equation*}
M_n={}_n(M_\ell\otimes D_\ell)
\end{equation*}
contains the image of the group ${}_nA(K')$ of $K'$-valued points of
${}_nA$. For $K'$ sufficiently large this group is equal to
${}_nA(\overline K)$, and therefore maps onto $M_n$. If
$\operatorname{char}K=p$, on the other hand, it appears that there is no
``trivial'' proof---that is, no proof avoiding \SGARefLink{sga7.sga7i.eix.n9-4}{9.4}---of \SGARefLink{sga7.sga7i.eix.n11-10}{11.10} for the
$p$-primary component.

b) The existence of an isomorphism

\begin{equation*}
\Psi\simeq(\mathbb Q/\mathbb Z)^r,
\end{equation*}
together with the correct value of $r$, was conjectured in 1964 by
J.-P. Serre.
